\section{Introduction}
\label{sec:intro}

\subsection{The problem}
\label{sec:problem}

Stateless LLM APIs require the full conversation to be sent at each turn. When
conversations exceed the context window (typically 128K--200K tokens), some form
of compaction is necessary. The naive approach --- summarize and discard old
messages --- is universally adopted by production systems (Claude Code, Cursor,
Windsurf, MemGPT/Letta). But how much information survives this process?

Answering this question requires a reliable measurement tool. Most existing
benchmarks treat recall measurement as straightforward: embed facts, ask
questions, count correct answers. We found that the measurement itself is
surprisingly fragile --- sensitive to how many questions are asked at once,
which categories of facts are tested, and whether ``present in context'' actually
means ``retrievable by the model.''

This led us to a three-phase approach: first calibrate the recall benchmark
(\S\ref{sec:calibration}), then use it to measure single-pass compaction loss
(\S\ref{sec:controlled}--\S\ref{sec:singlepass}), then compare multi-pass strategies
at scale (\S\ref{sec:strategy}). The calibration itself yielded findings
that challenge common assumptions about LLM recall.

\subsection{Related work}
\label{sec:related}

\textbf{Needle in a Haystack} \citep{kamradt2023needle}. The foundational test for long-context
retrieval: embed a specific fact in a long document and measure whether the
model can find it. Tests the model's \textit{native} retrieval ability within a single
context window. Our work extends this paradigm to measure fact \textit{survival}
across compaction cycles --- ``Needle in a Compacted Haystack.''

\textbf{LongMemEval} \citep{wu2024longmemeval}. Benchmarks long-term memory capabilities
of LLM systems across 500 questions in 6 categories (knowledge update, temporal
reasoning, multi-session, preferences, etc.) spanning up to 115 sessions. We
use LongMemEval evidence as our fact source, providing ecologically valid,
naturally-embedded facts rather than synthetic injections.

\textbf{Context Rot} \citep{chroma2025contextrot}. Demonstrates that LLM performance
degrades systematically as input length increases, even on trivially simple
tasks. Evaluates 18 models and recommends compaction/summarization as
mitigation, but does not compare compaction strategies against each other
or measure information loss across strategies.

\textbf{MemGPT / Letta} \citep{packer2023memgpt}. Proposes virtual context management
inspired by OS memory hierarchies: main context (RAM) and external context
(disk). When context overflows, messages are evicted and recursively
summarized --- functionally identical to our incremental strategy. The paper
notes that ``older messages have progressively less influence on the summary''
but does not benchmark this degradation or compare alternative strategies.

\textbf{Factory.ai --- Evaluating Context Compression} \citep{factory2025evaluating}. The closest
work to ours. Compares three compression \textit{implementations} (Factory's
anchored iterative, OpenAI Compact, Anthropic SDK) using probe-based
evaluation (recall, artifact, continuation, decision probes) on 36{,}000
production messages. A companion blog post on the design of Factory's own
anchored-iterative method \citep{factory2025compressing} is also cited where relevant. Key
differences with our work:
\begin{itemize}
  \item Factory compares \textit{implementations} from different vendors; we compare
    fundamentally different \textit{architectures} (single-shot vs iterative vs frozen)
  \item No zone-based recall metrics --- their evaluation does not reveal where
    in the conversation information is lost
  \item No frozen summary strategy or space-time tradeoff analysis
  \item No calibration of the recall measurement itself
\end{itemize}

\textbf{Beyond a Million Tokens} \citep{tavakoli2025beyondmillion}. Benchmarks long-term memory up to 10M
tokens across diverse domains, testing recall, multi-hop reasoning,
contradiction resolution, and temporal ordering. Tests model \textit{capabilities}
at various context lengths, not compaction \textit{strategy} effectiveness.

\textbf{Anthropic --- Effective Context Engineering} \citep{anthropic2025context}. Engineering guide
describing best practices for context management in AI agents, including
structured summarization and context pruning. Practical guidance but no
comparative benchmarking of strategies.

\textbf{Prompt Compression} \citep{jiang2023llmlingua}. Achieves up to
20x token compression with $\sim$1.5\% performance loss. Focuses on \textit{prompt-level}
compression (removing redundant tokens), orthogonal to \textit{conversation-level}
compaction (summarizing message history).

\subsection{Gap in existing work}
\label{sec:gap}

No prior work addresses the \textbf{measurement reliability} of compaction
benchmarks. Existing evaluations assume that embedding a fact and asking about
it provides a clean signal --- we show this assumption is wrong.

Beyond measurement, no prior work compares fundamentally different compaction
\textit{architectures} (single-shot vs iterative vs frozen summaries) on the same
conversation with spatial recall metrics. Existing benchmarks either test model
capabilities (Needle in a Haystack, Beyond a Million Tokens), measure context
rot without compaction (Chroma), or compare vendor implementations without
varying the underlying strategy (Factory.ai).

The re-summarization cascade problem (which we term the \textbf{``JPEG cascade''}:
each compaction cycle re-summarizes content that already includes earlier
summaries, accumulating loss like repeated lossy compression of an image)
is acknowledged implicitly in MemGPT but never quantified. The frozen
summary strategy and the resulting space-time tradeoff have not been
explored.

\subsection{Contributions}
\label{sec:contributions}

We introduce:

\begin{enumerate}
  \item \textbf{Evidence that recall measurement is non-trivial}: a significant questions-per-prompt
    effect (up to 20pp between Q=1 and Q=10) that \textbf{reverses direction} under severe
    compaction, a category-dependent recall hierarchy, and a systematic gap between
    keyword presence and LLM retrieval
  \item \textbf{A calibrated benchmark protocol} using LongMemEval evidence with
    controlled density, factorial evidence design, and repeatability validation
  \item \textbf{Controlled single-pass compaction loss measurement} isolating information
    destruction from context size reduction via re-padding
  \item \textbf{The attention dilution effect}: compaction degrades recall even in
    untouched portions of the context
  \item \textbf{The grep-LLM gap}: keywords survive compaction (82--93\% grep recall)
    but the LLM fails to retrieve them (0--7\% recall) --- a ``Lost in the Middle''
    amplification effect
  \item \textbf{Multi-pass strategy comparison} across conversation sizes (500K--10M
    tokens) showing that architectural optimization yields diminishing returns
    against the attention bottleneck
\end{enumerate}
